\section{Architecture}
\label{sec:arch}

\begin{figure}[t]
\centering
\begin{tikzpicture}[
  font=\scriptsize,
  box/.style={draw, rounded corners=1.5pt, align=center, minimum height=0.55cm, inner sep=2pt, fill=white},
  store/.style={draw, cylinder, shape border rotate=90, aspect=0.18, align=center, minimum height=0.62cm, inner sep=1.5pt, fill=white},
  arr/.style={-{Stealth[length=1.6mm]}, semithick},
  darr/.style={{Stealth[length=1.6mm]}-{Stealth[length=1.6mm]}, semithick},
  node distance=0.28cm and 0.3cm]

% Row 1: user-facing
\node[box, minimum width=2.0cm] (ui) {Desktop UI\\(React/Tauri)};
\node[box, right=of ui, minimum width=2.1cm] (api) {Task API\\(FastAPI, loopback)};
\node[box, right=of api, minimum width=1.7cm] (obs) {Observatory\\(separate process)};

% Row 2: runner + agent
\node[box, below=0.45cm of api, minimum width=5.6cm, minimum height=0.95cm] (agent) {};
\node[anchor=north, font=\scriptsize\bfseries] at (agent.north) {LangGraph agent (Fig.~\ref{fig:graph})};
\node[box, font=\tiny, anchor=south west, minimum width=1.25cm] (plan) at ($(agent.south west)+(0.08,0.07)$) {plan / ground};
\node[box, font=\tiny, right=0.08cm of plan, minimum width=0.95cm] (exec) {execute};
\node[box, font=\tiny, right=0.08cm of exec, minimum width=0.95cm] (ver) {verify $\Phi$};
\node[box, font=\tiny, right=0.08cm of ver, minimum width=1.05cm] (rec) {recover};

\node[box, left=0.25cm of agent, minimum width=1.2cm] (runner) {Task runner\\checkpoints};

% Row 3: services
\node[box, below=0.5cm of agent.south west, anchor=north west, minimum width=1.7cm] (router) {Model router};
\node[box, right=0.2cm of router, minimum width=1.75cm] (sec) {Sanitizer,\\approval gate};
\node[box, right=0.2cm of sec, minimum width=1.7cm] (pool) {Browser pool\\(Playwright)};

% Row 4: local resources
\node[box, below=0.35cm of router, minimum width=1.7cm, fill=gray!8] (ollama) {Ollama (local)\\text + vision};
\node[store, below=0.35cm of sec] (db) {SQLite\\tasks, events};
\node[store, right=0.12cm of db] (ev) {Evidence\\PNG, JSON};
\node[store, left=0.12cm of db] (chroma) {Chroma};
\node[box, below=0.35cm of pool, minimum width=1.7cm, fill=red!6, xshift=0.45cm] (web) {Web pages\\(untrusted)};

% Edges
\draw[darr] (ui) -- (api);
\draw[arr] (api) -- (agent.north -| api);
\draw[arr] (api.south west) ++(0,0) -- (runner.north);
\draw[arr] (runner) -- (agent);
\draw[darr] (router.north) -- (router.north |- agent.south);
\draw[darr] (router) -- (ollama);
\draw[arr] (sec.north) -- (sec.north |- agent.south);
\draw[darr] (pool.north) -- (pool.north |- agent.south);
\draw[darr] (pool) -- (web);
\draw[arr] (agent.south) ++(0.3,0) -- ++(0,-0.25) -| (ev.north);
\draw[arr, dashed] (db.east) ++(0,0.1) -- ++(0.12,0) |- (obs.south);
\draw[arr, dashed] (api.east) -- (obs.west);

% Trust boundary: everything except the web runs on the user's machine
\begin{scope}[on background layer]
\node[draw, dashed, rounded corners=3pt, inner sep=3pt, fit=(ui)(obs)(runner)(ollama)(chroma)(ev)(pool), label={[font=\tiny, anchor=south east, inner sep=1pt]south east:user's machine}] {};
\end{scope}
\end{tikzpicture}
\caption{Components of Agentic Pilot. Everything inside the dashed boundary runs on the user's machine; model inference goes to a local Ollama server. Web pages are the main untrusted input. Dashed arrows are the telemetry feed read by the Observatory.}
\label{fig:arch}
\end{figure}


Fig.~\ref{fig:arch} shows the components. A desktop client submits a natural-language task to a FastAPI service bound to the loopback interface. The task runner executes the control loop, checkpoints its state after every node and cancels tasks that exceed a wall-clock limit (5\,min by default). Task events go to SQLite and to a read-only monitoring process.

\subsection{Control Loop}
\label{sec:loop}

The loop is a compiled LangGraph state machine with eleven nodes; Algorithm~\ref{alg:loop} gives its behavior. A typed state record carries the intent, the plan and step index, the observation, the planned action, the action history, retry and model-call counters, routing decisions and the status. Only the verification step can end a task as \textsc{Completed}, and only when $\Phi$ holds (Fig.~\ref{fig:verify}).

\begin{algorithm}[t]
\caption{Evidence-gated control loop (browser backend)}
\label{alg:loop}
\footnotesize
\begin{algorithmic}[1]
\Require goal $g$; budget of $K{=}15$ model calls
\State $I \gets \Call{ParseIntent}{g}$; $P \gets \Call{Decompose}{g}$ if multi-step
\If{$I$ is high-risk and not approved} \Return \textsc{WaitApproval} \EndIf
\State \Call{Navigate}{$I.\mathit{site}$}; on failure go to line~\ref{alg:recover}
\Loop
  \State $o \gets \Call{Observe}{}$ \Comment{sanitized DOM, screenshot}
  \If{$\beta(o)$} \Return \textsc{Blocked} \Comment{hand off to user} \label{alg:beta}
  \EndIf
  \State $a \gets \Call{Plan}{g, P, o, \textsc{Retrieve}(g)}$ \Comment{model call}
  \State $\alpha \gets \Call{Execute}{a}$; save before/after evidence
  \If{$\alpha = 0$}
    \If{fewer than 3 retries} \textbf{continue} \EndIf
    \State $\textit{st} \gets \Call{Recover}{\text{error}}$ \label{alg:recover}
    \If{$\textit{st}$ is exhausted or blocked} \Return \textsc{Failed} \EndIf
    \State \textbf{continue} \Comment{retry, vision, \dots}
  \EndIf
  \If{$a \neq \texttt{complete}$} mark step done; \textbf{if} steps remain \textbf{continue} \EndIf
  \If{$a = \texttt{complete}$ or last step or text typed}
    \State $(v, \text{obs}) \gets \Phi(g, s, P)$; save as evidence \label{alg:phi}
    \If{$v$} store memory; \Return \textsc{Completed} \EndIf
  \EndIf
  \If{model calls $\geq K$} \Return \textsc{Failed} \EndIf
\EndLoop
\end{algorithmic}
\end{algorithm}


\paragraph{Planning} A small model parses the request into an intent (action, site, target, content, risk level). Requests with a connective or more than eight words are decomposed into steps with an action type and expected outcome. Each iteration, the planner sees the goal, the plan, the current observation and retrieved memory and returns one action. Deterministic rules then adjust it: exact-text requests are grounded to the best visible text field, repeated navigation becomes a failure, and plan steps that name clicking or extracting map to those actions.

\paragraph{Recovery, routing and memory} A failed action is retried through a fresh observation up to three times; later failures are classified into seven types and assigned a strategy (retry, alternative selector, vision fallback, replan), with at most six attempts. Only the vision-fallback strategy changes behavior, by routing the next step to the vision model; the other strategies repeat ordinary planning. A deterministic router assigns each model call a role and the first installed candidate among local Qwen2.5, Qwen3.5 and DeepSeek-R1 models (1.5B--7B parameters), or \texttt{qwen3-vl:2b} and \texttt{moondream} for vision. At termination, outcomes are written to a local Chroma collection and retrieved into later planning prompts.

\begin{figure}[t]
\centering
\begin{tikzpicture}[
  x=0.99cm, y=1.12cm,
  font=\scriptsize,
  b/.style={draw, rounded corners=1.5pt, align=center, inner sep=2pt, minimum height=0.5cm, fill=white},
  e/.style={draw, align=left, inner sep=1.5pt, font=\tiny, minimum width=2.4cm, fill=gray!6},
  c/.style={draw, circle, inner sep=0.5pt, minimum size=0.42cm, font=\tiny, fill=white},
  arr/.style={-{Stealth[length=1.4mm]}, semithick},
  lab/.style={font=\tiny, inner sep=1pt}]

\node[b] (act) at (0,0) {action $a_t$\\executed};
\node[b] (alpha) at (0,-1.15) {dispatch\\$\alpha_t\in\{0,1\}$};
\node[b] (rec) at (0,-2.45) {recovery: retry,\\vision, \dots};

\node[e] (ev1) at (2.35, 0.55) {URL, title, frames\\(error, sign-in, bot check)};
\node[e] (ev2) at (2.35,-0.15) {search parameters\\of the results page};
\node[e] (ev3) at (2.35,-0.85) {values of inputs\\and editable regions};
\node[e] (ev4) at (2.35,-1.55) {answer / structured\\extraction};
\node[e] (ev5) at (2.35,-2.25) {plan and step\\status};

\node[c] (cB) at (4.15, 0.75) {$\beta$};
\node[c] (cE) at (4.15, 0.25) {$E$};
\node[c] (cH) at (4.15,-0.25) {$H$};
\node[c] (cQ) at (4.15,-0.75) {$Q$};
\node[c] (cT) at (4.15,-1.25) {$T$};
\node[c] (cX) at (4.15,-1.75) {$X$};
\node[c] (cR) at (4.15,-2.25) {$R$};

\node[b, fill=gray!20, very thick] (and) at (5.35,-1.0) {$\Phi$\\$=\bigwedge$};
\node[b] (ok) at (6.95,-0.35) {\textsc{Completed}\\proof, memory};
\node[b] (no) at (6.95,-1.45) {reject: observe\\and replan};
\node[b] (blk) at (6.95,0.75) {\textsc{Blocked}: hand-off};
\node[draw, dashed, align=center, inner sep=2pt, font=\tiny] (store) at (6.95,-2.45) {evidence: decision +\\observed values};

\draw[arr] (act) -- (alpha);
\draw[arr] (alpha) -- node[lab, left] {0} (rec);
\draw[arr] (alpha.east) -- node[lab, below] {1} ++(0.35,0) |- (ev1.west);
\foreach \n in {ev2,ev3,ev4,ev5} \draw[arr] ($(alpha.east)+(0.35,0)$) |- (\n.west);
\draw[arr] (ev1.east) -- (cB.west); \draw[arr] (ev1.east) -- (cE.west); \draw[arr] (ev1.east) -- (cH.west);
\draw[arr] (ev2.east) -- (cQ.west); \draw[arr] (ev3.east) -- (cT.west);
\draw[arr] (ev4.east) -- (cX.west); \draw[arr] (ev5.east) -- (cR.west);
\foreach \n in {cE,cH,cQ,cT,cX,cR} \draw[arr] (\n.east) -- (and.west);
\draw[arr] (cB.east) -- (blk.west);
\draw[arr] (and.east) -- ++(0.25,0) |- node[lab, right, pos=0.3] {holds} (ok.west);
\draw[arr] (and.east) -- ++(0.25,0) |- node[lab, right, pos=0.3] {fails} (no.west);
\draw[arr, dashed] (and.south) |- (store.west);
\end{tikzpicture}
\caption{Evidence-gated verification after an action (Algorithm~\ref{alg:loop}, lines~\ref{alg:beta} and~\ref{alg:phi}). A failed dispatch goes to recovery and never to $\Phi$. A successful dispatch makes the observed values available to the conjuncts of the browser predicate, each triggered by the wording of the goal (Section~\ref{sec:problem}); the bot-check condition $\beta$ is tested first and hands the task to the user. Every decision is stored with the values it used.}
\label{fig:verify}
\end{figure}


\subsection{Browser Backend}
\label{sec:browser}

A pool of Playwright contexts over one Chromium instance gives each task its own page. Each observation is a DOM pass over interactable elements, recording tag, role, label, up to 120 characters of text, type, visibility, bounding box and selectors; elements are ranked, deduplicated and capped at 50, and their text is passed through a prompt-injection sanitizer. Actions are navigate, click, type, key press, scroll and extract; typed text is read back. The vision model is used only when the planner returns \texttt{need\_help} or recovery requests it: it receives a screenshot and returns a point, which is clicked without an individual check. Per task, the evidence directory holds before/after screenshots of each action, the latest DOM manifest, the verification and recovery records and the node trace.
